\section{Conclusion}
\label{sec:conclusion}

We evaluated source-grounded repair of document-field knowledge graphs using
structural normalization, diagnostic-context generation, and deterministic
candidate validation. The archived pipeline achieves 98.00\% controlled
defect repair and 99.17\% triple F1. On naturally extracted graphs, it improves
F1 from 87.99\% to 92.48\%; the gain remains after separating malformed JSON
inputs, although recovery of fully correct graphs remains limited.

A deterministic source-field parser reconstructs every reference graph,
because the current evidence explicitly serializes those fields. Thus these
benchmarks do not show that LLM repair is necessary or superior to direct
source reconstruction. Independently annotated unstructured-source data are
required to resolve this limitation.

The execution audit further narrows the interpretation of these results. The earlier
sequential optimizer does not explain the headline gains, and reference-aware
router features cannot support a deployment claim. A matched Gemma experiment
and an adapted SHACL-context baseline show no diagnostic-context advantage:
the base-plus-gate arm is better on controlled F1, while filtering yields only
a modest response-paired gain on natural inputs. The method remains a task-specific pipeline whose
literal-support and cardinality checks have limited semantic guarantees.
Independent human adjudication, broader graph structures, and validation of
relation-level source entailment remain necessary before making stronger
claims about general KG repair.
